\begin{table}[ht]
\caption{\textbf{\earth example}. The positive document describes the Airy's isostasy model, which is needed in calculating the hydrostatic equilibrium.}
\centering
% \resizebox{\textwidth}{!}{
\begin{tabular}{p{13cm}}
\toprule
% \hspace{5.5cm} Earth science \\
% \midrule
\textit{\textbf{Query}} \\
\midrule
How to calculate hydrostatic equilibrium?
\bigbreak
I'm trying to solve the following problem. The sea level in the past was 200 m higher than today. The seawater became in isostatic equilibrium with the ocean basin. what is the increase in the depth x of the ocean basins? Water density is $p_w = 1000$ kg $m^{-3}$, and mantle density is 3300 kg $m^{-3}$ \\
\bigbreak
Using the compensation column, I reach: \\
$x = (p_w * 200 m)$/3300 = 60.60 m \\
but normally I expected to find 290 m. \\
Can someone explain to me what's wrong? \\
\midrule
\textit{\textbf{Chain-of-thought reasoning to find documents}} \\
\midrule
This problem is about the hydrostatic equilibrium. 
This may be related to Pascal's law, which gives insights on the pressure change in fluid's mechanics.
To understand hydrostatic equilibrium or Pascal's law, it would be helpful to provide a detailed explanation of the concept and formula used in calculations.
Other content on the application of the theorems, or related lemmas that also focus on the discussion of isotatic equilibrium with the ocean basin would also be useful.\\
\midrule
\textit{\textbf{Example positive document}} \\
\midrule
Airy
\bigbreak
The basis of the model is Pascal's law, and particularly its consequence that, within a fluid in static equilibrium, the hydrostatic pressure is the same on every point at the same elevation (surface of hydrostatic compensation): \\
$h_1 \cdot p_1 = h_2 \cdot p_2 = h_3 \cdot p_3 = ... h_n \cdot p_n$ \\
For the simplified picture shown, the depth of the mountain belt roots (b1) is calculated as follows: \\
$(h_1+c+b_1)p_c = (cp_c)+(b_1p_m)$ \\
$b_1(p_m-p_c)=h_1p_c$ \\
$b_1 = h_1p_c/(p_m-p_c)$ \\
where $p_m$ is the density of the mantle (ca. 3,300 kg $m^{-3}$) and 
$p_c$ is the density of the crust (ca. 2,750 kg $m^{-3}$). Thus, generally: \\
$b_1 = 5 h_1$ \\
... \\
\midrule
\textit{\textbf{Example negative document}} \\
\midrule
Global or eustatic sea level has fluctuated significantly over Earth's history. The main factors affecting sea level are the amount and volume of available water and the shape and volume of the ocean basins. The primary influences on water volume are the temperature of the seawater, which affects density, and the amounts of water retained in other reservoirs like rivers, aquifers, lakes, glaciers, polar ice caps and sea ice. Over geological timescales, changes in the shape of the oceanic basins and in land/sea distribution affect sea level. In addition to eustatic changes, local changes in sea level are caused by the earth's crust uplift and subsidence. \\

Over geologic time sea level has fluctuated by more than 300 metres, possibly more than 400 metres. The main reasons for sea level fluctuations in the last 15 million years are the Antarctic ice sheet and Antarctic post-glacial rebound during warm periods. \\

The current sea level is about 130 metres higher than the historical minimum. Historically low levels were reached during the Last Glacial Maximum (LGM), about 20,000 years ago. The last time the sea level was higher than today was during the Eemian, about 130,000 years ago. \\
... \\
\bottomrule
\end{tabular}
% }
\label{tab:earth_science_example}
\end{table}